\documentclass[10pt,a4paper]{amsart}
\usepackage[margin=19mm]{geometry}
\usepackage{amsmath,amssymb,mathtools}
\usepackage[scr=rsfs,cal=euler]{mathalfa}
\usepackage{xfrac}
\usepackage[unicode,hidelinks,bookmarks=false]{hyperref}
\newcommand{\ie}{i.e.\ }
\newcommand{\cf}{cf.\ }
\newcommand{\resp}{respectively\ }
\newcommand{\Subsection}{\subsection}
\providecommand{\nobreakdash}{\nobreak\hbox{-}}
\setlength{\emergencystretch}{2em}
\multlinegap0pt
\usepackage{dsfont}
\usepackage{amssymb}
\usepackage{mathtools}
\newtheorem{proposition}{Proposition}
\newtheorem{lemma}{Lemma}
\newcommand{\CF}{\mathcal F}
\newcommand{\CL}{{\mathcal L}{}}
\newcommand{\G}{\mathcal{G}}
\newcommand\Op{\operatorname{Op}}
\newcommand{\vf}{\varphi}
\title{Quantization of locally compact groups associated with essentially bijective \texorpdfstring{$1$}{1}-cocycles}
\author{Pierre Bieliavsky, Victor Gayral, Sergey Neshveyev, Lars Tuset}
\date{Source: arXiv:2312.00523v1 (CC BY 4.0)}
\begin{document}
\maketitle

\noindent\emph{Source and licence.} Quantization of locally compact groups associated with essentially bijective \texorpdfstring{$1$}{1}-cocycles. Version arXiv:2312.00523v1 (CC BY 4.0), distributed by arXiv under the Creative Commons Attribution 4.0 International licence, http://creativecommons.org/licenses/by/4.0/, as recorded in the arXiv licence metadata for this exact version and shown on the abstract page https://arxiv.org/abs/2312.00523v1. Full licence notice: \texttt{licenses/arxiv-2312.00523-CC-BY-4.0.txt}.\par\smallskip

\noindent\textit{Editorial scope.} Counted unit: the article's lemma Gomega (source0.tex lines 868--884). The article's complete original proof of it is delivered with it (source0.tex lines 886--975, one \texttt{proof} environment of 90 lines). No mathematics is written, repaired or re-derived: the statement and the proof are the article's own lines, in the article's own notation, and no retained line is substituted or re-flowed.  Retained uncounted as the source-local context the proof needs: lines 412--420; lines 558--561; lines 861--863.  Nothing was omitted from any retained span, so the package carries no omission record.  No retained line contains a citation, so the wrapper carries no bibliography and no bibliography entry.  No retained line contains a figure environment, an image-inclusion command or drawing code, so no image is delivered and the package declares no asset dependency.  Editorial formatting only, in addition: an AMS \texttt{amsart} preamble that restates the article's own declarations and macros used by the retained lines, namely the environments \texttt{proposition}, \texttt{lemma} on separate counters, together with the standard packages those lines need.  The retained environments print in this document as Proposition 1, Lemma 1 in order of appearance, whereas the article prints the same items with the numbering of its own full document.  The complete native source of the article as distributed in the arXiv e-print for this exact version is bundled unchanged as \texttt{sources/150824/source0.tex}.
\begin{proposition}\label{prop:representation}
Define a map $\pi\colon G\to U(L^2(Q))$ by
\begin{align}
\label{def-pi}
(\pi(q,v)\vf)(q_0):=e^{-i \langle\eta(q_0),v\rangle}\,b (q_0^{-1},q)\,\vf\big(q^{-1}q_0\big),\quad \vf\in L^2(Q).
\end{align}
This is a projective unitary representation with cocycle $\omega$, so that $\pi(g_1)\pi(g_2)=\omega(g_1,g_2)\,\pi(g_1g_2)$ for all $g_1,g_2\in G$.
Furthermore, $\pi$ is irreducible, square-integrable and the corresponding Duflo--Moore operator is the operator of multiplication by the function $\Delta^{-1}|_Q=|\cdot|\Delta_Q^{-1}$.
\end{proposition}

\begin{align}
 \label{Komega}
 K(f)(q_0,q)=|q_0q|^{-1/2}\,b (q^{-1}_0,q_0)\,\overline{b (q^{-1},q_0)}\,(\CF_V f)(q_0, \eta(q_0)-\eta(q)).
\end{align}

\begin{equation}\label{eq:FL-invariance}
(\CF_V\lambda_{g} f)(q',\xi')=|q| e^{-i\langle \xi',v+\beta(q,q^{-1}q')\rangle}(\CF_Vf)(q^{-1}q',{q^{-1}}^\flat\xi').
\end{equation}

\begin{lemma}
\label{Gomega}
For every $f\in L^2(Q\times\hat V\times Q\times\hat V,|q|^{-1}dq\times d\xi\times |q|^{-1}dq\times d\xi)$, we have:

\medskip
$\displaystyle ((\CF_V\otimes \CF_V)\tilde\G(\CF_V^*\otimes \CF_V^*)f)(q_1,\xi_1;q_2,\xi_2)$
\begin{align*}
=\ &
|q_1|\Delta(q_1)^{-1/2}\Delta(\eta^{-1}(\xi_1+\eta(q_2)))^{1/2}\,e^{i\langle\xi_1,\beta(q_1,q_1^{-1}q_2)\rangle} \\
& b (\eta^{-1}(\eta(q_2)-\xi_2)^{-1},q_2)\,
\overline{b (\eta^{-1}(\xi_1+\eta(q_2))^{-1},q_2)}\\
& b (\eta^{-1}(\xi_1+\eta(q_2))^{-1},\eta^{-1}(\xi_1+\eta(q_2)))\,\overline{b(\eta^{-1}(\eta(q_2)-\xi_2)^{-1},\eta^{-1}(\xi_1+\eta(q_2)))}
\\
& f(q_1^{-1}q_2,-{q_1^{-1}}^\flat\xi_1;
\eta^{-1}(\xi_1+\eta(q_2)), \xi_1+\xi_2).
\end{align*}
\end{lemma}

\begin{proof}
If $f\in \CF\CL_0(G)$ and $\CF_Vf$ is zero outside $X_{L,M}$, then the function $K(f)$ given by~\eqref{Komega} is zero outside $L\times M$. It follows that if $D$ is the Duflo--Moore operator from Proposition~\ref{prop:representation}, then $\Op(f)D^{-1/2}$ is a kernel operator with kernel $K(f)(1\otimes
\Delta^{1/2}|_Q)$, which is a bounded function vanishing outside $L\times M$. In particular, $\Op(f)D^{-1/2}$ is  Hilbert--Schmidt and thus bounded.

Take $f_1,f_2\in \CF\CL_0(G)$. Then for the Galois map $\tilde \G$ we have
\begin{align*}
(\tilde \G(f_1\otimes f_2))(g_1,g_2)&=\Op^*\big(\pi (g_1)\Op(f_1)D^{-1/2}\pi (g_1)^*
\Op(f_2)\big)(g_2)\\
&=\Delta(g_1)^{-1/2}\Op^*\big(\Op(\lambda_{g_1}f_1)D^{-1/2}
\Op(f_2)\big)(g_2).
\end{align*}
The operator $\Op(\lambda_{g_1}f_1)D^{-1/2}\Op(f_2)$ has kernel $ k_{g_1}\in L^2(Q\times Q)$,
given by
\begin{align*}
k_{g_1}(q_2,q_3)
&= \int K(\lambda_{g_1}f_1)(q_2,q_0)\,\Delta(q_0)^{1/2} \,
 K(f_2)(q_0,q_3)\,dq_0\\
 &= \int b (q_2^{-1},q_2)
\,\overline{b(q_0^{-1},q_2)}\,b (q_0^{-1},q_0)
\,\overline{b(q_3^{-1},q_0)}\,\Delta(q_0)^{1/2}
 \\
& \qquad(\CF_V \lambda_{g_1}f_1)(q_2, \eta(q_2)-\eta(q_0))\,(\CF_V f_2)(q_0, \eta(q_0)-\eta(q_3))
 \,|q_2q_3|^{-1/2}\,|q_0|^{-1}\,dq_0.
 \end{align*}

We need to write this as $K(f)$ for some $f$. For this we have to invert the kernel formula~\eqref{Komega}. Letting
$$
k(q_0,q):=|q_0q|^{-1/2}\,b (q_0^{-1},q_0)
\,\overline{b (q^{-1},q_0)}\,(\CF_V f)(q_0, \eta(q_0)-\eta(q))
$$
and $\xi= \eta(q_0)-\eta(q)$, so that $q=\eta^{-1}(\eta(q_0)-\xi)$, we get
\begin{multline*}
(\CF_V f)(q_0, \xi)\\=|q_0\eta^{-1}(\eta(q_0)-\xi)|^{-1/2}\,\overline{b (q_0^{-1},q_0)}
\,{b }(\eta^{-1}(\eta(q_0)-\xi)^{-1},q_0)\,k(q_0,\eta^{-1}(\eta(q_0)-\xi)).
\end{multline*}

We therefore deduce:

\medskip
$\displaystyle \big((1\otimes \CF_V)\tilde \G(f_1\otimes f_2)\big)(q_1,v_1;q_2,\xi_2)$
\begin{align*}
&=\Delta(q_1)^{-1/2}\,
|q_2\eta^{-1}(\eta(q_2)-\xi_2)|^{1/2}\,\overline{b(q_2^{-1},q_ 2)}\,b (\eta^{-1}(\eta(q_2)-\xi_2)^{-1},q_2)\\
&\qquad
k_{g_1}(q_2,\eta^{-1}(\eta(q_2)-\xi_2))\\
&=\Delta(q_1)^{-1/2}
\, b (\eta^{-1}(\eta(q_2)-\xi_2)^{-1},q_2)\int \overline{b(q_0^{-1},q_2)}\,b (q_0^{-1},q_0)
\,\overline{b(\eta^{-1}(\eta(q_2)-\xi_2)^{-1},q_0)} \\
& \qquad(\CF_V \lambda_{g_1}f_1)(q_2, \eta(q_2)-\eta(q_0))\,(\CF_V f_2)(q_0, \eta(q_0)-\eta(q_2)+\xi_2)
 \,\Delta(q_0)^{1/2}\,|q_0|^{-1}\,dq_0.
\end{align*}
Applying~\eqref{eq:FL-invariance} we then get

\medskip
$\displaystyle \big((1\otimes \CF_V)\tilde \G(f_1\otimes f_2)\big)(q_1,v_1;q_2,\xi_2)$
\begin{align}\label{eq:star-product0}
=\ &|q_1|\,\Delta(q_1)^{-1/2}
\, b (\eta^{-1}(\eta(q_2)-\xi_2)^{-1},q_2)\notag\\
&\int \overline{b (q_0^{-1},q_2)}\,b (q_0^{-1},q_0)
\,\overline{b (\eta^{-1}(\eta(q_2)-\xi_2)^{-1},q_0)}\notag \\
&(\CF_Vf_1)(q_1^{-1}q_2,{q_1^{-1}}^\flat(\eta(q_2)-\eta(q_0)))\,(\CF_V f_2)(q_0, \eta(q_0)-\eta(q_2)+\xi_2)\notag\\
&e^{-i\langle \eta(q_2)-\eta(q_0),v_1+\beta(q_1,q_1^{-1}q_2)\rangle}\,\Delta(q_0)^{1/2}\,\frac{dq_0}{\,|q_0|}.
\end{align}
Setting $\xi_1:=\eta(q_0)$, we get

\medskip
$\displaystyle \big((1\otimes \CF_V)\tilde \G(f_1\otimes f_2)\big)(q_1,v_1;q_2,\xi_2)$
\begin{align*}
=\ &|q_1|\,\Delta(q_1)^{-1/2}
\, b (\eta^{-1}(\eta(q_2)-\xi_2)^{-1},q_2)\\
&\int \overline{b (\eta^{-1}(\xi_1)^{-1},q_2)}\,b (\eta^{-1}(\xi_1)^{-1},\eta^{-1}(\xi_1))
\,\overline{b (\eta^{-1}(\eta(q_2)-\xi_2)^{-1},\eta^{-1}(\xi_1))}\\
& (\CF_Vf_1)(q_1^{-1}q_2,{q_1^{-1}}^\flat(\eta(q_2)-\xi_1))\,(\CF_V f_2)(\eta^{-1}(\xi_1),\xi_1-\eta(q_2)+\xi_2)\\
& e^{-i\langle \eta(q_2)-\xi_1,v_1+\beta(q_1,q_1^{-1}q_2)\rangle}\,\Delta(\eta^{-1}(\xi_1))^{1/2}\,d\xi_1.
\end{align*}
Performing the translation $\xi_1\mapsto \xi_1+\eta(q_2)$ and using ${q^{-1}}^\flat\eta^{-1}(\xi+\eta(q))=\eta^{-1}( {q^{-1}}^\flat\xi)$, we get

\medskip
$\displaystyle \big((1\otimes \CF_V)\tilde \G(f_1\otimes f_2)\big)(q_1,v_1;q_2,\xi_2)$
\begin{align*}
=\ &|q_1|\,\Delta(q_1)^{-1/2}
\, b (\eta^{-1}(\eta(q_2)-\xi_2)^{-1},q_2)\\
&\int \overline{b (\eta^{-1}(\xi_1+\eta(q_2))^{-1},q_2)}\,b (\eta^{-1}(\xi_1+\eta(q_2))^{-1},\eta^{-1}(\xi_1+\eta(q_2)))\\
&\overline{b (\eta^{-1}(\eta(q_2)-\xi_2)^{-1},\eta^{-1}(\xi_1+\eta(q_2)))}\\
& (\CF_Vf_1)(q_1^{-1}q_2,-{q_1^{-1}}^\flat\xi_1)\,(\CF_V f_2)(\eta^{-1}(\xi_1+\eta(q_2)),\xi_1+\xi_2)\\
& e^{i\langle\xi_1,v_1+\beta(q_1,q_1^{-1}q_2)\rangle}\,\Delta(\eta^{-1}(\xi_1+\eta(q_2)))^{1/2}\,d\xi_1.
\end{align*}

Applying $\CF_V\otimes1$ we conclude that the lemma is true for the functions $f=\CF_Vf_1\otimes \CF_Vf_2$, with $f_1,f_2\in \CF\CL_0(G)$, which is enough by density of the linear span of such functions.
\end{proof}

\end{document}
